\documentclass[10pt,a4paper]{amsart}
\usepackage{amsmath,amssymb,amsthm,graphicx}
% Retained minimal preamble for the Stacks Project chapter more-morphisms.tex (More on Morphisms).
% Verbatim excerpt of the chapter preamble: Xy-pic setup (native preamble lines 14-18) and the native
% theorem declarations with the mathematical macros (native preamble lines 154-217). Copied mechanically;
% nothing added, removed or redefined. The wrapper class and the amsmath/amssymb/amsthm/graphicx base
% packages are supplied by the assembler, so the class/externaldocument part of the native preamble is omitted.
% For commutative diagrams we use Xy-pic
\usepackage[all]{xy}

% We use 2cell for 2-commutative diagrams.
\xyoption{2cell}
% Theorem environments.
%
\theoremstyle{plain}
\newtheorem{theorem}[subsection]{Theorem}
\newtheorem{proposition}[subsection]{Proposition}
\newtheorem{lemma}[subsection]{Lemma}

\theoremstyle{definition}
\newtheorem{definition}[subsection]{Definition}
\newtheorem{example}[subsection]{Example}
\newtheorem{exercise}[subsection]{Exercise}
\newtheorem{situation}[subsection]{Situation}

\theoremstyle{remark}
\newtheorem{remark}[subsection]{Remark}
\newtheorem{remarks}[subsection]{Remarks}

\numberwithin{equation}{subsection}

% Macros
%
\def\lim{\mathop{\mathrm{lim}}\nolimits}
\def\colim{\mathop{\mathrm{colim}}\nolimits}
\def\Spec{\mathop{\mathrm{Spec}}}
\def\Hom{\mathop{\mathrm{Hom}}\nolimits}
\def\Ext{\mathop{\mathrm{Ext}}\nolimits}
\def\SheafHom{\mathop{\mathcal{H}\!\mathit{om}}\nolimits}
\def\SheafExt{\mathop{\mathcal{E}\!\mathit{xt}}\nolimits}
\def\Sch{\mathit{Sch}}
\def\Mor{\mathop{\mathrm{Mor}}\nolimits}
\def\Ob{\mathop{\mathrm{Ob}}\nolimits}
\def\Sh{\mathop{\mathit{Sh}}\nolimits}
\def\NL{\mathop{N\!L}\nolimits}
\def\CH{\mathop{\mathrm{CH}}\nolimits}
\def\proetale{{pro\text{-}\acute{e}tale}}
\def\etale{{\acute{e}tale}}
\def\QCoh{\mathit{QCoh}}
\def\Ker{\mathop{\mathrm{Ker}}}
\def\Im{\mathop{\mathrm{Im}}}
\def\Coker{\mathop{\mathrm{Coker}}}
\def\Coim{\mathop{\mathrm{Coim}}}

% Boxtimes
%
\DeclareMathSymbol{\boxtimes}{\mathbin}{AMSa}{"02}

%
% Macros for moduli stacks/spaces
%
\def\QCohstack{\mathcal{QC}\!\mathit{oh}}
\def\Cohstack{\mathcal{C}\!\mathit{oh}}
\def\Spacesstack{\mathcal{S}\!\mathit{paces}}
\def\Quotfunctor{\mathrm{Quot}}
\def\Hilbfunctor{\mathrm{Hilb}}
\def\Curvesstack{\mathcal{C}\!\mathit{urves}}
\def\Polarizedstack{\mathcal{P}\!\mathit{olarized}}
\def\Complexesstack{\mathcal{C}\!\mathit{omplexes}}
% \Pic is the operator that assigns to X its picard group, usage \Pic(X)
% \Picardstack_{X/B} denotes the Picard stack of X over B
% \Picardfunctor_{X/B} denotes the Picard functor of X over B
\def\Pic{\mathop{\mathrm{Pic}}\nolimits}
\def\Picardstack{\mathcal{P}\!\mathit{ic}}
\def\Picardfunctor{\mathrm{Pic}}
\def\Deformationcategory{\mathcal{D}\!\mathit{ef}}
\newenvironment{reference}{}{}
\newenvironment{slogan}{}{}
\newenvironment{history}{}{}
\usepackage{cleveref}
\crefname{originallemma}{original lemma}{original lemmas}
\Crefname{originallemma}{Original lemma}{Original lemmas}
\makeatletter
\newlabel{algebra-lemma-colimit-faithfully-flat}{{algebra-lemma-colimit-faithfully-flat}{source}{}{sourcepointer-1}{}}
\newlabel{algebra-lemma-colimit-faithfully-flat@cref}{{[originallemma][0][]algebra-lemma-colimit-faithfully-flat}{[1][0][]source}{}{}{}}
\makeatother
\crefname{originallemma}{original lemma}{original lemmas}
\Crefname{originallemma}{Original lemma}{Original lemmas}
\makeatletter
\newlabel{algebra-lemma-ff-rings}{{algebra-lemma-ff-rings}{source}{}{sourcepointer-2}{}}
\newlabel{algebra-lemma-ff-rings@cref}{{[originallemma][0][]algebra-lemma-ff-rings}{[1][0][]source}{}{}{}}
\makeatother
\crefname{originallemma}{original lemma}{original lemmas}
\Crefname{originallemma}{Original lemma}{Original lemmas}
\makeatletter
\newlabel{algebra-lemma-flat-going-down}{{algebra-lemma-flat-going-down}{source}{}{sourcepointer-3}{}}
\newlabel{algebra-lemma-flat-going-down@cref}{{[originallemma][0][]algebra-lemma-flat-going-down}{[1][0][]source}{}{}{}}
\makeatother
\crefname{originallemma}{original lemma}{original lemmas}
\Crefname{originallemma}{Original lemma}{Original lemmas}
\makeatletter
\newlabel{algebra-lemma-minimal-prime-image-minimal-prime}{{algebra-lemma-minimal-prime-image-minimal-prime}{source}{}{sourcepointer-4}{}}
\newlabel{algebra-lemma-minimal-prime-image-minimal-prime@cref}{{[originallemma][0][]algebra-lemma-minimal-prime-image-minimal-prime}{[1][0][]source}{}{}{}}
\makeatother
\title{The number of minimal primes of a filtered colimit of faithfully flat maps is the supremum of the stages}
\author{The Stacks Project Authors}
\date{stacks-project:more-morphisms.tex:local-snapshot-20261011; GNU FDL 1.2 or later}
\begin{document}
\maketitle
\noindent\textit{Editorial note.} Statement span is lines 10302-10309 of sources/stacks-more-morphisms/pinned-original-source.tex; proof span is lines 10311-10331, the complete author proof (19 body lines).\par\medskip
\noindent\textit{Editorial note.} Navigation stubs: 4 cited-but-not-included labels are supplied as non-load-bearing pointers with their pinned tags where the snapshot tag index has them.\par\medskip
\noindent\textit{Editorial note.} extra\_packages is empty on purpose: the retained preamble already loads xy with options and a second declaration clashes in the assembler wrapper.\par\medskip
\noindent\textit{Editorial note.} Text is retained verbatim, including the author's own typos and placeholders.\par\medskip
\noindent\textit{Allowed standard prerequisites (editorial).}\begin{itemize}
\item Explicit prerequisite boundary. Minimal primes of a ring, filtered colimits of rings, and faithful flatness; the comparison is retained.
\end{itemize}\medskip
\ifdefined\hypertarget\hypertarget{sourcepointer-1}{}\fi
\noindent\textit{Source reference.} Original lemma algebra-lemma-colimit-faithfully-flat: sources/stacks-more-morphisms/pinned-original-source.tex, source lines 10331--10331. Cited by the retained statement or proof but not part of the collected spans; supplied only as a navigation pointer. The argument does not depend on its content.\par\medskip
\ifdefined\hypertarget\hypertarget{sourcepointer-2}{}\fi
\noindent\textit{Source reference.} Original lemma algebra-lemma-ff-rings: sources/stacks-more-morphisms/pinned-original-source.tex, source lines 10331--10331. Cited by the retained statement or proof but not part of the collected spans; supplied only as a navigation pointer. The argument does not depend on its content.\par\medskip
\ifdefined\hypertarget\hypertarget{sourcepointer-3}{}\fi
\noindent\textit{Source reference.} Original lemma algebra-lemma-flat-going-down: sources/stacks-more-morphisms/pinned-original-source.tex, source lines 10331--10331. Cited by the retained statement or proof but not part of the collected spans; supplied only as a navigation pointer. The argument does not depend on its content.\par\medskip
\ifdefined\hypertarget\hypertarget{sourcepointer-4}{}\fi
\noindent\textit{Source reference.} Original lemma algebra-lemma-minimal-prime-image-minimal-prime: sources/stacks-more-morphisms/pinned-original-source.tex, source lines 10331--10331. Cited by the retained statement or proof but not part of the collected spans; supplied only as a navigation pointer. The argument does not depend on its content.\par\medskip
\begin{lemma}
\label{lemma-nr-minimal-primes}
Let $R = \colim R_i$ be colimit of a directed system of rings
whose transition maps are faithfully flat.
Then the number of minimal primes of $R$
taken as an element of $\{0, 1, 2, \ldots, \infty\}$
is the supremum of the numbers of minimal primes of the $R_i$.
\end{lemma}

\begin{proof}
If $A \to B$ is a flat ring map, then $\Spec(B) \to \Spec(A)$
maps minimal primes to minimal primes by going down
(Algebra, Lemma \ref{algebra-lemma-flat-going-down}).
If $A \to B$ is faithfully flat, then every minimal
prime is the image of a minimal prime (by
Algebra, Lemma \ref{algebra-lemma-ff-rings} and
\ref{algebra-lemma-minimal-prime-image-minimal-prime}).
Hence the number of minimal primes of $R_i$ is
$\geq$ the number of minimal primes of $R_{i'}$ if $i \leq i'$.
By Algebra, Lemma \ref{algebra-lemma-colimit-faithfully-flat}
each of the maps $R_i \to R$ is
faithfully flat and we also see that
the number of minimal primes of $R$ is
$\geq$ the number of minimal primes of $R_i$.
Finally, suppose that $\mathfrak q_1, \ldots, \mathfrak q_n$
are pairwise distinct minimal primes of $R$. Then we can
find an $i$ such that $R_i \cap \mathfrak q_1, \ldots, R_i \cap \mathfrak q_n$
are pairwise distinct (as sets and hence as prime ideals).
This implies the lemma.
\end{proof}

\end{document}
